% 辛群（综述）
% license CCBY4
% type Wiki

本文根据 CC-BY-SA 协议转载翻译自维基百科\href{https://en.wikipedia.org/wiki/Symplectic_group}{相关文章}

在数学中，“辛群”的名称可以指两类不同但密切相关的数学群族，分别记作Sp(2n, F)与Sp(n)，其中n正整数，F为域（通常取C或 R）。后一种被称为紧致辛群，也记作$\mathrm{USp}(n)$。许多作者使用略有不同的记号，通常差别体现在系数 2 的处理上。此处采用的记号与表示这些群的最常见矩阵的规模保持一致。在卡当（Cartan）对单李代数的分类中，复群Sp(2n, C) 的李代数记作Cₙ，而Sp(n)则是Sp(2n, C)的紧致实型。需要注意的是，当我们提及“（紧致）辛群”时，通常是指依照维数n编号的一族（紧致）辛群。

“辛群”这一名称由赫尔曼·外尔提出，用来取代此前令人困惑的名称——“（线）复群”（(line) complex group）和“阿贝尔线性群”。它是“复数”一词的希腊语类比。

梅塔辛群是实数域上辛群的双覆盖；它在其他局部域、有限域以及阿德尔环上也有类似的对应物。
\subsection{Sp(2n, F)}
辛群是一个经典群，定义为作用在2n 维向量空间（定义在域$F$上）上的线性变换的集合，这些变换保持一个非退化的斜对称双线性型。这样的向量空间称为辛向量空间，而一个抽象辛向量空间$V$的辛群记作$\mathrm{Sp}(V)$。一旦为$V$选定基底，辛群就变为由矩阵乘法构成的$2n \times 2n$辛矩阵群，其元素属于域$F$。这个群记作$\mathrm{Sp}(2n, F)$或$\mathrm{Sp}(n, F)$。如果该双线性型由一个非奇异斜对称矩阵 $\Omega$ 表示，则有：
$$
\operatorname{Sp}(2n,F) = \{ M \in M_{2n \times 2n}(F) : M^{\mathrm{T}} \, \Omega \, M = \Omega \},~
$$
其中$M^{\mathrm{T}}$表示矩阵$M$的转置。通常，$\Omega$定义为：
$$
\Omega = 
\begin{pmatrix}
0 & I_n \\
- I_n & 0
\end{pmatrix},~
$$
其中$I_n$是 $n \times n$ 单位矩阵。在这种情况下，$\mathrm{Sp}(2n, F)$ 可以表示为如下的分块矩阵：
$$
\begin{pmatrix}
A & B \\
C & D
\end{pmatrix}, \quad A, B, C, D \in M_{n \times n}(F),~
$$
并满足以下三个方程：
$$
\begin{aligned}
- C^{\mathrm{T}} A + A^{\mathrm{T}} C &= 0, \\
- C^{\mathrm{T}} B + A^{\mathrm{T}} D &= I_n, \\
- D^{\mathrm{T}} B + B^{\mathrm{T}} D &= 0.
\end{aligned}~
$$
